\section{Introduction}
\label{sec:intro}

A classical physics-informed neural network (PINN) learns low-frequency solution
components fast and high-frequency components almost not at all --- spectral bias --- and
its neural tangent kernel (NTK) eigenspectrum decays with Fourier-mode index, so residual
and boundary loss terms converge at characteristically different rates. Separately, a
data-re-uploading variational quantum circuit is known to realize a truncated Fourier
series whose accessible frequency set is exactly determined by its data-encoding
Hamiltonians' eigenvalues \citep{schuld2021effect}. These two facts have not previously
been joined into a constructive design rule: prior hybrid-QPINN work
\citep{hybrid_qpinn_mlst,hybrid_when_where_2606,qpinn_maxwell_2506,cq_hybrid_arch_2511}
reports empirical accuracy deltas on hand-picked circuit sizes, without a principled way
to choose the circuit \emph{given} the PDE.

This paper makes that connection. Given a PDE, we compute its target spectral support
analytically from the operator's Fourier symbol and forcing/boundary data, then construct
the quantum circuit's encoding generators, re-uploading depth, and qubit count so its
realized frequency set covers that support --- no architecture search, no hyperparameter
sweep over circuit size. We call this Spectrum-Matched Circuit Design (SMCD,
Section~\ref{sec:methodology}). We then use six explainable-AI instruments, run identically for
matched classical and hybrid models at the same training checkpoints, as the measurement
instrument that attributes any accuracy delta to the quantum layer specifically ---
not as after-the-fact decoration on a result already claimed by other means.

\paragraph{Contributions.}
\begin{description}
\item[C1 --- SMCD, a constructive PDE $\to$ circuit map.] Given a linear operator's symbol
  and forcing/boundary spectra, output encoding-generator eigenvalues, re-uploading depth,
  qubit count, and entangler pattern, deterministically. Falsifiable: a spectrum-matched
  circuit must beat a size-matched random-frequency circuit on the same PDE, or the design
  step is not earning its keep.
\item[C2 --- Quantum-NTK block decomposition.] $\Theta_{\mathrm{hyb}} = \Theta_{\mathrm{cl}}
  + \Theta_q$ for disjoint-parameter hybrid models (Section~\ref{sec:prop4}), with the quantum
  block's eigenvalue decay flattened over the encoded band relative to the classical
  block's power-law decay.
\item[C3 --- Spectral-bias relief as mechanism, not just outcome.] Per-frequency error
  trajectories must show the hybrid model's high-$k$ modes converging at rates comparable
  to its low-$k$ modes, with the improvement band coinciding with the encoded frequency
  set --- not merely an aggregate accuracy improvement.
\item[C4 --- A negative-result map with a decision rule.] An explicit characterization of
  PDE classes where hybridization is useless or harmful, with the mathematical reason
  (smoothing operators remove high-frequency content, so a high-frequency-rich circuit
  adds parameters and barren-plateau risk for no benefit).
\item[C5 --- An XAI protocol for hybrid physics-informed models.] A reusable,
  model-agnostic instrument set with published metric definitions, run identically across
  classical and quantum families rather than as a post hoc, model-specific explanation.
\end{description}

The honest headline this paper reports, whatever the specific numbers turn out to be
(Section~\ref{sec:results}--\ref{sec:negative}): the quantum layer should help on high-wavenumber
and multiscale problems where its encoded band matches the PDE's spectral support, do
nothing on smooth diffusive problems whose target spectrum SMCD predicts to be
low-frequency-only \emph{before any training}, and hurt when the circuit's spectrum is
mismatched to the problem. A design rule plus a stated ``when not to go quantum''
recommendation, adjudicated against pre-registered falsifiers (Section~\ref{sec:results}), is
the contribution --- not a cherry-picked win.

\paragraph{Assignment mapping (T5.9 self-review).}

\begin{center}
\begin{tabular}{p{0.42\textwidth}p{0.5\textwidth}}
\toprule
\textbf{Requirement} & \textbf{Where it is satisfied} \\
\midrule
Solve a set of PDEs with classical PINNs and QAPINNs &
  Section~\ref{sec:methodology} (4 PDE instances $\times$ 6 problem labels, \S\ref{sec:algorithm}),
  Section~\ref{sec:methodology} (7 model families), one training harness for both \\
XAI techniques/metrics showing \emph{how} learning happens in each &
  Section~\ref{sec:xai} (6-instrument suite), run on both families at matched checkpoints,
  results in Section~\ref{sec:results} \\
Demonstrate how the VQC affects the classical PINN learning process &
  Section~\ref{sec:results} (NTK block decay, PR-7; per-frequency convergence, PR-8; matched-
  capacity ablation, PR-1--3), Section~\ref{sec:negative} (negative controls, PR-4, PR-11) \\
Methodology to construct a problem-specific circuit + architecture, with mathematical basis &
  Section~\ref{sec:methodology} (SMCD, Propositions 1--4, Algorithm 1) \\
Technical report, 5--15 pp &
  This document (15 pp, \texttt{paper/main.pdf}) \\
Source code repository &
  \texttt{src/qapinn/}, \texttt{README.md} \\
Reproducibility instructions &
  \texttt{docs/REPRODUCE.md}, \texttt{README.md}, \texttt{python tasks.py repro-quick} \\
Presentation slides &
  \texttt{slides/main.pdf}, 18 slides, same figure assets as this document \\
Summary of key findings + recommendations &
  \texttt{FINDINGS.md} (13 numbered findings, C1--C5 table, decision table) \\
\bottomrule
\end{tabular}
\end{center}

Self-review against this table (T5.9): every figure cited above appears in
Section~\ref{sec:results}--\ref{sec:negative} with the specific prediction it decides; no
verdict stated anywhere in this paper exceeds the corresponding entry in
\texttt{docs/predictions.md}'s mechanical adjudication
(\texttt{scripts/adjudicate\_predictions.py}) --- Section~\ref{sec:conclusion}'s C1--C5
summary is a synthesis of those verdicts, not an independent claim.
